\subsection{Modelling Random Experiments}

Probability theory is about modelling and analysing \textit{random experiments}: experiments whose outcome cannot be determined in advance, but are nevertheless still subject to analysis. Mathematically, we can model a random experiment by defining a specific measure space \((\Omega, \cH, \Pm)\), which we call a \textit{probability space}, where the three components are: a \textit{sample space}, a collection of \textit{events}, and a \textit{probability measure}.

The \textit{sample space} \(\Omega\) of a random experiment is the set of all possible outcomes of the random experiment. Sample spaces can be countable, such as \(\N\), or uncountable, such as \(\R^n\) or \(\{0, 1\}^\infty\). The sets in the \(\sigma\)-algebra \(\cH\) are called \textit{events}. These are the subsets to which we wish to assign a probability. We can view \(\cH\) as the family of events. An event \(A\) \textit{occurs} if the outcome of the experiment is one of the elements in \(A\). The properties of \(\cH\) are natural in the context of random experiments and events:

\begin{itemize}
	\item If \(A\) and \(B\) are events, the set \(A \cup B\) is also an event, namely the event that \(A\) \textit{or} \(B\) or both occur. For a sequence \(A_1, A_2, \dots\) of events, their union is the event that \textit{at least one} of the events occurs.
	\item If \(A\) and \(B\) are events, the set \(A \cap B\) is also an event, namely the event that \(A\) \textit{and} \(B\) both occur. For a sequence \(A_1, A_2, \dots\) of events, their intersection is the event that \textit{all} of the events occur.
	\item If \(A\) is an event, its complement \(A^c\) is also an event, namely the event that \(A\) \textit{does not occur}.
	\item The set \(\Omega\) itself is an event, namely the \textit{certain event} (it always occurs). Similarly, \(\emptyset\) is an event, namely the \textit{impossible event} (it never occurs).
\end{itemize}

When the sample space \(\Omega\) is \(\R\), the natural \(\sigma\)-algebra \(\cH\) is the \textit{Borel \(\sigma\)-algebra} \(\cB\). Recall that this is the smallest \(\sigma\)-algebra on \(\R\) that contains all the intervals of the form \((-\infty, x]\) for \(x \in \R\). This \(\sigma\)-algebra of sets is big enough to contain all important sets and small enough to allow us to assign a probability to all events.

The third ingredient in the model for a random experiment is the specification of the probability of the events. In fact, this is the crucial part of the model. Mathematically, we are looking for a \textit{measure} \(\Pm\) that assigns to each event \(A\) a number \(\Pm(A)\) in \([0, 1]\) describing how likely or probable that event is.

\begin{defn}{Probability Measure}{ProbabilityMeasure}
	A \textit{probability measure} is a measure on \((\Omega, \cH)\) with \(\Pm(\Omega) = 1\). Thus, \(\Pm\) is a mapping from \(\cH\) to \([0, 1]\) with the following properties:
	\begin{enumerate}
		\item \(\Pm(\emptyset) = 0\) and \(\Pm(\Omega) = 1\).
		\item For any sequence \(A_1, A_2, \dots\) of disjoint events,
		      \begin{equation}
			      \Pm(\cup_n A_n) = \sum_n \Pm(A_n).
			      \label{eq:ProbabilityMeasureDef}
		      \end{equation}
	\end{enumerate}
\end{defn}

An event \(A\) is said to occur \textit{almost surely} (a.s.) if \(\Pm(A) = 1\).

\begin{thm}{Properties of a Probability Measure}{PropertiesOfAProbabilityMeasure}
	Let \(\Pm\) be a probability measure on a measurable space \((\Omega, \cH)\), and let \(A, B\) and \(A_1, A_2, \dots\) be events. Then, the following hold:
	\begin{enumerate}
		\item (\textit{Monotonicity}): \(A \subseteq B \implies \Pm(A) \le \Pm(B)\).
		\item (\textit{Continuity from below}): If \(A_1, A_2, \dots\) is an increasing sequence of events, i.e., \(A_1 \subseteq A_2 \subseteq \cdots\), then
		      \begin{equation}
			      \lim \Pm(A_n) = \Pm(\cup_n A_n).
			      \label{eq:ContinuityFromBelow}
		      \end{equation}
		\item (\textit{Countable subadditivity}): \(\Pm(\cup_n A_n) \le \sum_n \Pm(A_n)\).
	\end{enumerate}
\end{thm}

\begin{thm}{Additional Properties of a Probability Measure}{}
	Let \(\Pm\) be a probability measure on a measurable space \((\Omega, \cH)\), and let \(A, B\) and \(A_1, A_2, \dots\) be events. Then, the following hold:
	\begin{enumerate}
		\item (\textit{Complement}): \(\Pm(A^c) = 1 - \Pm(A)\).
		\item (\textit{Union}): \(\Pm(A \cup B) = \Pm(A) + \Pm(B) - \Pm(A \cap B)\).
		\item (\textit{Continuity from above}): If \(A_1, A_2, \dots\) is a decreasing sequence of events, i.e., \(A_1 \supseteq A_2 \supseteq \cdots\), then
		      \begin{equation}
			      \lim \Pm(A_n) = \Pm(\cap_n A_n).
			      \label{eq:ContinuityFromAbove}
		      \end{equation}
	\end{enumerate}
\end{thm}

\subsection{Random Variables}
Usually, the most convenient way to describe quantities of interest connected with random experiments is by using random variables. Random variables allow us to use intuitive notation for certain events, such as \(\{X \in A\}\), etc. Intuitively, we can think of a random variable \(X\) as a measurement on a random experiment that will become available \textit{tomorrow}. However, all the thinking work can be done \textit{today}. If for each possible outcome \(\omega\), the measurement is \(X(\omega)\), we can specify today probabilities such as \(\Pm(X \in A)\).

We can translate our intuitive notion of a random variable into rigorous mathematics by defining a random variable to be a measurable function from \((\Omega, \cH)\) to a measurable space \((E, \cE)\).

\begin{defn}{Random Variable}{RandomVariable}
	Let \((\Omega, \cH, \Pm)\) be a probability space and \((E, \cE)\) a measurable space. A \textit{random variable} \(X\) taking values in \(E\) is an \(\cH/\cE\)-measurable function; that is, a mapping \(X : \Omega \to E\) that satisfies
	\begin{equation}
		X^{-1}A = \{X \in A\} = \{\omega \in \Omega : X(\omega) \in A\} \in \cH \quad \text{for all } A \in \cE.
		\label{eq:RandomVariableDef}
	\end{equation}
\end{defn}

Often a random experiment is described via more than one random variable. Let \((\Omega,\cH,\Pm)\) be a probability space and let \(\mathbb{T}\) be an arbitrary index set, which may be countable or uncountable. Think of \(\mathbb{T}\) as a set of times. For each \(t \in \mathbb{T}\), let \((E_t,\cE_t)\) be a measurable space. The product measurable space of \(\{(E_t,\cE_t): t \in \mathbb{T}\}\) is
\[
	\Big(\times_{t \in \mathbb{T}} E_t,\;\bigotimes_{t \in \mathbb{T}} \cE_t\Big).
\]
Each element \(x \in \times_{t \in \mathbb{T}} E_t\) is of the form \(x = (x_t)_{t \in \mathbb{T}}\). Often \((E_t,\cE_t)\) is the same space \((E,\cE)\) for all \(t\), in which case \(x\) can be viewed as a function \(x:\mathbb{T} \to E\). The \(\sigma\)-algebra on \(\times_{t \in \mathbb{T}} E_t\) is generated by the rectangles
\begin{equation}
	R\bigl((A_t)_{t \in \mathbb{T}}\bigr) \idef \Bigl\{\,x \in \times_{t \in \mathbb{T}} E_t : x_t \in A_t \text{ for each } t \in \mathbb{T}\,\Bigr\},
	\label{eq:ProductRectangles}
\end{equation}
where \(A_t \in \cE_t\) and \(A_t = E_t\) except for finitely many \(t\).

\begin{defn}{Stochastic Process}{}
	A collection of random variables \(X \coloneqq \{X_t, t \in \mathbb{T}\}\), where \(\mathbb{T}\) is any index set, is called a \textit{stochastic process}. When \(\mathbb{T}\) is finite, \(X\) is called a \textit{random vector}.
\end{defn}

\begin{thm}{Stochastic Process}{}
	\(X \coloneqq \{X_t, t \in \mathbb{T}\}\) is a stochastic process taking values in \(\times_{t \in \mathbb{T}} E_t\) with \(\sigma\)-algebra \(\otimes_{t \in \mathbb{T}} \cE_t\) if and only if each \(X_t\) is a random variable taking values in \(E_t\) with \(\sigma\)-algebra \(\cE_t\).
\end{thm}

\subsection{Probability Distributions}
Let \((\Omega, \cH, \Pm)\) be a probability space and let \(X\) be a random variable taking values
in a set \(E\) with \(\sigma\)-algebra \(\cE\). To simplify our usage, we also say that ``\(X\) takes values
in \((E, \cE)\)''. If we wish to describe our experiment via \(X\), then we need to specify
the probabilities of events such as \(\{X \in B\}, B \in \cE\). In elementary probability
theory we usually made a distinction between ``discrete'' and ``continuous'' random
variables. However, it is sometimes better to analyse random variables in a more
unified manner; many definitions and properties of random variables do not depend
on whether they are ``discrete'' or ``continuous''. Moreover, some random variables
are neither discrete nor continuous.

Recall that by definition \textit{all} the sets \(\{X \in B\}\) with \(B \in \cE\) are events. So we
may assign a probability to each of these events. This leads to the notion of the
\textit{distribution} of a random variable.

\begin{defn}{Probability Distribution}{}
	Let \((\Omega, \cH, \Pm)\) be a probability space and let \(X\) be a random variable taking values in \((E, \cE)\). The (\textit{probability}) \textit{distribution} of \(X\) is the image measure \(\mu\) of \(\Pm\) under \(X\). That is,
	\begin{equation}
		\mu(B) \coloneqq \Pm(X \in B), \quad B \in \cE.
		\label{eq:ProbabilityDistributionDef}
	\end{equation}
\end{defn}

\begin{defn}{Cumulative Distribution Function}{}
	The (\textit{cumulative}) \textit{distribution function} of a numerical random variable \(X\) is the function \(F : \R \to [0, 1]\) defined by
	\begin{equation}
		F(x) \coloneqq \Pm(X \le x), \quad x \in \R.
		\label{eq:CumulativeDistributionFunctionDef}
	\end{equation}
\end{defn}

\begin{prop}{Properties of a Cumulative Distribution Function}{}
	\begin{enumerate}
		\item (\textit{Bounded}): \(0 \le F(x) \le 1\).
		\item (\textit{Increasing}): \(x \le y \implies F(x) \le F(y)\).
		\item (\textit{Right-continuous}): \(\lim_{h \downarrow 0} F(x + h) = F(x)\).
	\end{enumerate}
\end{prop}

\begin{defn}{Discrete Distribution}{}
	We say that \(X\) has a \textit{discrete distribution} \(\mu\) on \((E, \cE)\) if \(\mu\) is a discrete measure; that is, it is of the form
	\begin{equation*}
		\mu = \sum_{x \in D} f(x) \delta_x
	\end{equation*}
	for some countable set \(D\) and positive masses \(\{f(x), x \in D\}\).
\end{defn}

\begin{defn}{Distribution with a Density}{}
	We say that \(X\) has an \textit{absolutely continuous} distribution \(\mu\) on \((E, \cE)\) with respect to a measure \(\lambda\) on \((E, \cE)\) if there exists a positive \(\cE\)-measurable function \(f\) such that \(\mu = f \lambda\); that is,
	\begin{equation*}
		\mu(B) = \Pm(X \in B) = \int_B \lambda(\di x) f(x), \quad B \in \cE.
	\end{equation*}
\end{defn}

The function \(f\) is thus the \textit{(probability) density function} (pdf) of \(X\) with respect to the measure, \(\lambda\) often taken to be the Lebesgue measure. Note that \(\mu\) is absolutely continuous with respect to \(\lambda\) in the sense that for all \(B \in \cE: \lambda(B) = 0 \implies \mu(B) = 0\).
\begin{figure}[H]
	\centering
	%% Creator: Matplotlib, PGF backend
%%
%% To include the figure in your LaTeX document, write
%%   \input{<filename>.pgf}
%%
%% Make sure the required packages are loaded in your preamble
%%   \usepackage{pgf}
%%
%% Also ensure that all the required font packages are loaded; for instance,
%% the lmodern package is sometimes necessary when using math font.
%%   \usepackage{lmodern}
%%
%% Figures using additional raster images can only be included by \input if
%% they are in the same directory as the main LaTeX file. For loading figures
%% from other directories you can use the `import` package
%%   \usepackage{import}
%%
%% and then include the figures with
%%   \import{<path to file>}{<filename>.pgf}
%%
%% Matplotlib used the following preamble
%%   \def\mathdefault#1{#1}
%%   \everymath=\expandafter{\the\everymath\displaystyle}
%%   \IfFileExists{scrextend.sty}{
%%     \usepackage[fontsize=10.000000pt]{scrextend}
%%   }{
%%     \renewcommand{\normalsize}{\fontsize{10.000000}{12.000000}\selectfont}
%%     \normalsize
%%   }
%%   \usepackage{amsmath}
%%   \usepackage{amssymb}
%%   \usepackage{amsfonts}
%%   \usepackage{libertine}
%%   \usepackage[libertine]{newtxmath}
%%   \makeatletter\@ifpackageloaded{underscore}{}{\usepackage[strings]{underscore}}\makeatother
%%
\begingroup%
\makeatletter%
\begin{pgfpicture}%
\pgfpathrectangle{\pgfpointorigin}{\pgfqpoint{6.300000in}{4.025000in}}%
\pgfusepath{use as bounding box, clip}%
\begin{pgfscope}%
\pgfsetbuttcap%
\pgfsetmiterjoin%
\pgfsetlinewidth{0.000000pt}%
\definecolor{currentstroke}{rgb}{0.000000,0.000000,0.000000}%
\pgfsetstrokecolor{currentstroke}%
\pgfsetstrokeopacity{0.000000}%
\pgfsetdash{}{0pt}%
\pgfpathmoveto{\pgfqpoint{0.000000in}{0.000000in}}%
\pgfpathlineto{\pgfqpoint{6.300000in}{0.000000in}}%
\pgfpathlineto{\pgfqpoint{6.300000in}{4.025000in}}%
\pgfpathlineto{\pgfqpoint{0.000000in}{4.025000in}}%
\pgfpathlineto{\pgfqpoint{0.000000in}{0.000000in}}%
\pgfpathclose%
\pgfusepath{}%
\end{pgfscope}%
\begin{pgfscope}%
\pgfpathrectangle{\pgfqpoint{0.050000in}{0.050000in}}{\pgfqpoint{6.200000in}{3.925000in}}%
\pgfusepath{clip}%
\pgfsetbuttcap%
\pgfsetroundjoin%
\definecolor{currentfill}{rgb}{0.678431,0.847059,0.901961}%
\pgfsetfillcolor{currentfill}%
\pgfsetfillopacity{0.900000}%
\pgfsetlinewidth{0.000000pt}%
\definecolor{currentstroke}{rgb}{0.000000,0.000000,0.000000}%
\pgfsetstrokecolor{currentstroke}%
\pgfsetdash{}{0pt}%
\pgfpathmoveto{\pgfqpoint{2.777628in}{3.539603in}}%
\pgfpathlineto{\pgfqpoint{2.777628in}{0.540625in}}%
\pgfpathlineto{\pgfqpoint{2.783682in}{0.540625in}}%
\pgfpathlineto{\pgfqpoint{2.789737in}{0.540625in}}%
\pgfpathlineto{\pgfqpoint{2.795792in}{0.540625in}}%
\pgfpathlineto{\pgfqpoint{2.801847in}{0.540625in}}%
\pgfpathlineto{\pgfqpoint{2.807902in}{0.540625in}}%
\pgfpathlineto{\pgfqpoint{2.813957in}{0.540625in}}%
\pgfpathlineto{\pgfqpoint{2.820011in}{0.540625in}}%
\pgfpathlineto{\pgfqpoint{2.826066in}{0.540625in}}%
\pgfpathlineto{\pgfqpoint{2.832121in}{0.540625in}}%
\pgfpathlineto{\pgfqpoint{2.838176in}{0.540625in}}%
\pgfpathlineto{\pgfqpoint{2.844231in}{0.540625in}}%
\pgfpathlineto{\pgfqpoint{2.850286in}{0.540625in}}%
\pgfpathlineto{\pgfqpoint{2.856340in}{0.540625in}}%
\pgfpathlineto{\pgfqpoint{2.862395in}{0.540625in}}%
\pgfpathlineto{\pgfqpoint{2.868450in}{0.540625in}}%
\pgfpathlineto{\pgfqpoint{2.874505in}{0.540625in}}%
\pgfpathlineto{\pgfqpoint{2.880560in}{0.540625in}}%
\pgfpathlineto{\pgfqpoint{2.886615in}{0.540625in}}%
\pgfpathlineto{\pgfqpoint{2.892669in}{0.540625in}}%
\pgfpathlineto{\pgfqpoint{2.898724in}{0.540625in}}%
\pgfpathlineto{\pgfqpoint{2.904779in}{0.540625in}}%
\pgfpathlineto{\pgfqpoint{2.910834in}{0.540625in}}%
\pgfpathlineto{\pgfqpoint{2.916889in}{0.540625in}}%
\pgfpathlineto{\pgfqpoint{2.922944in}{0.540625in}}%
\pgfpathlineto{\pgfqpoint{2.928999in}{0.540625in}}%
\pgfpathlineto{\pgfqpoint{2.935053in}{0.540625in}}%
\pgfpathlineto{\pgfqpoint{2.941108in}{0.540625in}}%
\pgfpathlineto{\pgfqpoint{2.947163in}{0.540625in}}%
\pgfpathlineto{\pgfqpoint{2.953218in}{0.540625in}}%
\pgfpathlineto{\pgfqpoint{2.959273in}{0.540625in}}%
\pgfpathlineto{\pgfqpoint{2.965328in}{0.540625in}}%
\pgfpathlineto{\pgfqpoint{2.971382in}{0.540625in}}%
\pgfpathlineto{\pgfqpoint{2.977437in}{0.540625in}}%
\pgfpathlineto{\pgfqpoint{2.983492in}{0.540625in}}%
\pgfpathlineto{\pgfqpoint{2.989547in}{0.540625in}}%
\pgfpathlineto{\pgfqpoint{2.995602in}{0.540625in}}%
\pgfpathlineto{\pgfqpoint{3.001657in}{0.540625in}}%
\pgfpathlineto{\pgfqpoint{3.007711in}{0.540625in}}%
\pgfpathlineto{\pgfqpoint{3.013766in}{0.540625in}}%
\pgfpathlineto{\pgfqpoint{3.019821in}{0.540625in}}%
\pgfpathlineto{\pgfqpoint{3.025876in}{0.540625in}}%
\pgfpathlineto{\pgfqpoint{3.031931in}{0.540625in}}%
\pgfpathlineto{\pgfqpoint{3.037986in}{0.540625in}}%
\pgfpathlineto{\pgfqpoint{3.044040in}{0.540625in}}%
\pgfpathlineto{\pgfqpoint{3.050095in}{0.540625in}}%
\pgfpathlineto{\pgfqpoint{3.056150in}{0.540625in}}%
\pgfpathlineto{\pgfqpoint{3.062205in}{0.540625in}}%
\pgfpathlineto{\pgfqpoint{3.068260in}{0.540625in}}%
\pgfpathlineto{\pgfqpoint{3.074315in}{0.540625in}}%
\pgfpathlineto{\pgfqpoint{3.080369in}{0.540625in}}%
\pgfpathlineto{\pgfqpoint{3.086424in}{0.540625in}}%
\pgfpathlineto{\pgfqpoint{3.092479in}{0.540625in}}%
\pgfpathlineto{\pgfqpoint{3.098534in}{0.540625in}}%
\pgfpathlineto{\pgfqpoint{3.104589in}{0.540625in}}%
\pgfpathlineto{\pgfqpoint{3.110644in}{0.540625in}}%
\pgfpathlineto{\pgfqpoint{3.116698in}{0.540625in}}%
\pgfpathlineto{\pgfqpoint{3.122753in}{0.540625in}}%
\pgfpathlineto{\pgfqpoint{3.128808in}{0.540625in}}%
\pgfpathlineto{\pgfqpoint{3.134863in}{0.540625in}}%
\pgfpathlineto{\pgfqpoint{3.140918in}{0.540625in}}%
\pgfpathlineto{\pgfqpoint{3.146973in}{0.540625in}}%
\pgfpathlineto{\pgfqpoint{3.153027in}{0.540625in}}%
\pgfpathlineto{\pgfqpoint{3.159082in}{0.540625in}}%
\pgfpathlineto{\pgfqpoint{3.165137in}{0.540625in}}%
\pgfpathlineto{\pgfqpoint{3.171192in}{0.540625in}}%
\pgfpathlineto{\pgfqpoint{3.177247in}{0.540625in}}%
\pgfpathlineto{\pgfqpoint{3.183302in}{0.540625in}}%
\pgfpathlineto{\pgfqpoint{3.189356in}{0.540625in}}%
\pgfpathlineto{\pgfqpoint{3.195411in}{0.540625in}}%
\pgfpathlineto{\pgfqpoint{3.201466in}{0.540625in}}%
\pgfpathlineto{\pgfqpoint{3.207521in}{0.540625in}}%
\pgfpathlineto{\pgfqpoint{3.213576in}{0.540625in}}%
\pgfpathlineto{\pgfqpoint{3.219631in}{0.540625in}}%
\pgfpathlineto{\pgfqpoint{3.225685in}{0.540625in}}%
\pgfpathlineto{\pgfqpoint{3.231740in}{0.540625in}}%
\pgfpathlineto{\pgfqpoint{3.237795in}{0.540625in}}%
\pgfpathlineto{\pgfqpoint{3.243850in}{0.540625in}}%
\pgfpathlineto{\pgfqpoint{3.249905in}{0.540625in}}%
\pgfpathlineto{\pgfqpoint{3.255960in}{0.540625in}}%
\pgfpathlineto{\pgfqpoint{3.262014in}{0.540625in}}%
\pgfpathlineto{\pgfqpoint{3.268069in}{0.540625in}}%
\pgfpathlineto{\pgfqpoint{3.274124in}{0.540625in}}%
\pgfpathlineto{\pgfqpoint{3.280179in}{0.540625in}}%
\pgfpathlineto{\pgfqpoint{3.286234in}{0.540625in}}%
\pgfpathlineto{\pgfqpoint{3.292289in}{0.540625in}}%
\pgfpathlineto{\pgfqpoint{3.298343in}{0.540625in}}%
\pgfpathlineto{\pgfqpoint{3.304398in}{0.540625in}}%
\pgfpathlineto{\pgfqpoint{3.310453in}{0.540625in}}%
\pgfpathlineto{\pgfqpoint{3.316508in}{0.540625in}}%
\pgfpathlineto{\pgfqpoint{3.322563in}{0.540625in}}%
\pgfpathlineto{\pgfqpoint{3.328618in}{0.540625in}}%
\pgfpathlineto{\pgfqpoint{3.334672in}{0.540625in}}%
\pgfpathlineto{\pgfqpoint{3.340727in}{0.540625in}}%
\pgfpathlineto{\pgfqpoint{3.346782in}{0.540625in}}%
\pgfpathlineto{\pgfqpoint{3.352837in}{0.540625in}}%
\pgfpathlineto{\pgfqpoint{3.358892in}{0.540625in}}%
\pgfpathlineto{\pgfqpoint{3.364947in}{0.540625in}}%
\pgfpathlineto{\pgfqpoint{3.371001in}{0.540625in}}%
\pgfpathlineto{\pgfqpoint{3.377056in}{0.540625in}}%
\pgfpathlineto{\pgfqpoint{3.383111in}{0.540625in}}%
\pgfpathlineto{\pgfqpoint{3.389166in}{0.540625in}}%
\pgfpathlineto{\pgfqpoint{3.395221in}{0.540625in}}%
\pgfpathlineto{\pgfqpoint{3.401276in}{0.540625in}}%
\pgfpathlineto{\pgfqpoint{3.407331in}{0.540625in}}%
\pgfpathlineto{\pgfqpoint{3.413385in}{0.540625in}}%
\pgfpathlineto{\pgfqpoint{3.419440in}{0.540625in}}%
\pgfpathlineto{\pgfqpoint{3.425495in}{0.540625in}}%
\pgfpathlineto{\pgfqpoint{3.431550in}{0.540625in}}%
\pgfpathlineto{\pgfqpoint{3.437605in}{0.540625in}}%
\pgfpathlineto{\pgfqpoint{3.443660in}{0.540625in}}%
\pgfpathlineto{\pgfqpoint{3.449714in}{0.540625in}}%
\pgfpathlineto{\pgfqpoint{3.455769in}{0.540625in}}%
\pgfpathlineto{\pgfqpoint{3.461824in}{0.540625in}}%
\pgfpathlineto{\pgfqpoint{3.467879in}{0.540625in}}%
\pgfpathlineto{\pgfqpoint{3.473934in}{0.540625in}}%
\pgfpathlineto{\pgfqpoint{3.479989in}{0.540625in}}%
\pgfpathlineto{\pgfqpoint{3.486043in}{0.540625in}}%
\pgfpathlineto{\pgfqpoint{3.492098in}{0.540625in}}%
\pgfpathlineto{\pgfqpoint{3.498153in}{0.540625in}}%
\pgfpathlineto{\pgfqpoint{3.504208in}{0.540625in}}%
\pgfpathlineto{\pgfqpoint{3.510263in}{0.540625in}}%
\pgfpathlineto{\pgfqpoint{3.516318in}{0.540625in}}%
\pgfpathlineto{\pgfqpoint{3.522372in}{0.540625in}}%
\pgfpathlineto{\pgfqpoint{3.528427in}{0.540625in}}%
\pgfpathlineto{\pgfqpoint{3.534482in}{0.540625in}}%
\pgfpathlineto{\pgfqpoint{3.540537in}{0.540625in}}%
\pgfpathlineto{\pgfqpoint{3.546592in}{0.540625in}}%
\pgfpathlineto{\pgfqpoint{3.552647in}{0.540625in}}%
\pgfpathlineto{\pgfqpoint{3.558701in}{0.540625in}}%
\pgfpathlineto{\pgfqpoint{3.564756in}{0.540625in}}%
\pgfpathlineto{\pgfqpoint{3.570811in}{0.540625in}}%
\pgfpathlineto{\pgfqpoint{3.576866in}{0.540625in}}%
\pgfpathlineto{\pgfqpoint{3.582921in}{0.540625in}}%
\pgfpathlineto{\pgfqpoint{3.588976in}{0.540625in}}%
\pgfpathlineto{\pgfqpoint{3.595030in}{0.540625in}}%
\pgfpathlineto{\pgfqpoint{3.601085in}{0.540625in}}%
\pgfpathlineto{\pgfqpoint{3.607140in}{0.540625in}}%
\pgfpathlineto{\pgfqpoint{3.613195in}{0.540625in}}%
\pgfpathlineto{\pgfqpoint{3.619250in}{0.540625in}}%
\pgfpathlineto{\pgfqpoint{3.625305in}{0.540625in}}%
\pgfpathlineto{\pgfqpoint{3.631359in}{0.540625in}}%
\pgfpathlineto{\pgfqpoint{3.637414in}{0.540625in}}%
\pgfpathlineto{\pgfqpoint{3.643469in}{0.540625in}}%
\pgfpathlineto{\pgfqpoint{3.649524in}{0.540625in}}%
\pgfpathlineto{\pgfqpoint{3.655579in}{0.540625in}}%
\pgfpathlineto{\pgfqpoint{3.661634in}{0.540625in}}%
\pgfpathlineto{\pgfqpoint{3.667688in}{0.540625in}}%
\pgfpathlineto{\pgfqpoint{3.673743in}{0.540625in}}%
\pgfpathlineto{\pgfqpoint{3.679798in}{0.540625in}}%
\pgfpathlineto{\pgfqpoint{3.685853in}{0.540625in}}%
\pgfpathlineto{\pgfqpoint{3.691908in}{0.540625in}}%
\pgfpathlineto{\pgfqpoint{3.697963in}{0.540625in}}%
\pgfpathlineto{\pgfqpoint{3.704017in}{0.540625in}}%
\pgfpathlineto{\pgfqpoint{3.710072in}{0.540625in}}%
\pgfpathlineto{\pgfqpoint{3.716127in}{0.540625in}}%
\pgfpathlineto{\pgfqpoint{3.722182in}{0.540625in}}%
\pgfpathlineto{\pgfqpoint{3.728237in}{0.540625in}}%
\pgfpathlineto{\pgfqpoint{3.734292in}{0.540625in}}%
\pgfpathlineto{\pgfqpoint{3.740346in}{0.540625in}}%
\pgfpathlineto{\pgfqpoint{3.746401in}{0.540625in}}%
\pgfpathlineto{\pgfqpoint{3.752456in}{0.540625in}}%
\pgfpathlineto{\pgfqpoint{3.758511in}{0.540625in}}%
\pgfpathlineto{\pgfqpoint{3.764566in}{0.540625in}}%
\pgfpathlineto{\pgfqpoint{3.770621in}{0.540625in}}%
\pgfpathlineto{\pgfqpoint{3.776675in}{0.540625in}}%
\pgfpathlineto{\pgfqpoint{3.782730in}{0.540625in}}%
\pgfpathlineto{\pgfqpoint{3.788785in}{0.540625in}}%
\pgfpathlineto{\pgfqpoint{3.794840in}{0.540625in}}%
\pgfpathlineto{\pgfqpoint{3.800895in}{0.540625in}}%
\pgfpathlineto{\pgfqpoint{3.806950in}{0.540625in}}%
\pgfpathlineto{\pgfqpoint{3.813004in}{0.540625in}}%
\pgfpathlineto{\pgfqpoint{3.819059in}{0.540625in}}%
\pgfpathlineto{\pgfqpoint{3.825114in}{0.540625in}}%
\pgfpathlineto{\pgfqpoint{3.831169in}{0.540625in}}%
\pgfpathlineto{\pgfqpoint{3.837224in}{0.540625in}}%
\pgfpathlineto{\pgfqpoint{3.843279in}{0.540625in}}%
\pgfpathlineto{\pgfqpoint{3.849333in}{0.540625in}}%
\pgfpathlineto{\pgfqpoint{3.855388in}{0.540625in}}%
\pgfpathlineto{\pgfqpoint{3.861443in}{0.540625in}}%
\pgfpathlineto{\pgfqpoint{3.867498in}{0.540625in}}%
\pgfpathlineto{\pgfqpoint{3.873553in}{0.540625in}}%
\pgfpathlineto{\pgfqpoint{3.879608in}{0.540625in}}%
\pgfpathlineto{\pgfqpoint{3.885662in}{0.540625in}}%
\pgfpathlineto{\pgfqpoint{3.891717in}{0.540625in}}%
\pgfpathlineto{\pgfqpoint{3.897772in}{0.540625in}}%
\pgfpathlineto{\pgfqpoint{3.903827in}{0.540625in}}%
\pgfpathlineto{\pgfqpoint{3.909882in}{0.540625in}}%
\pgfpathlineto{\pgfqpoint{3.915937in}{0.540625in}}%
\pgfpathlineto{\pgfqpoint{3.921992in}{0.540625in}}%
\pgfpathlineto{\pgfqpoint{3.928046in}{0.540625in}}%
\pgfpathlineto{\pgfqpoint{3.934101in}{0.540625in}}%
\pgfpathlineto{\pgfqpoint{3.940156in}{0.540625in}}%
\pgfpathlineto{\pgfqpoint{3.946211in}{0.540625in}}%
\pgfpathlineto{\pgfqpoint{3.952266in}{0.540625in}}%
\pgfpathlineto{\pgfqpoint{3.958321in}{0.540625in}}%
\pgfpathlineto{\pgfqpoint{3.964375in}{0.540625in}}%
\pgfpathlineto{\pgfqpoint{3.970430in}{0.540625in}}%
\pgfpathlineto{\pgfqpoint{3.976485in}{0.540625in}}%
\pgfpathlineto{\pgfqpoint{3.982540in}{0.540625in}}%
\pgfpathlineto{\pgfqpoint{3.988595in}{0.540625in}}%
\pgfpathlineto{\pgfqpoint{3.994650in}{0.540625in}}%
\pgfpathlineto{\pgfqpoint{4.000704in}{0.540625in}}%
\pgfpathlineto{\pgfqpoint{4.006759in}{0.540625in}}%
\pgfpathlineto{\pgfqpoint{4.012814in}{0.540625in}}%
\pgfpathlineto{\pgfqpoint{4.018869in}{0.540625in}}%
\pgfpathlineto{\pgfqpoint{4.024924in}{0.540625in}}%
\pgfpathlineto{\pgfqpoint{4.030979in}{0.540625in}}%
\pgfpathlineto{\pgfqpoint{4.037033in}{0.540625in}}%
\pgfpathlineto{\pgfqpoint{4.043088in}{0.540625in}}%
\pgfpathlineto{\pgfqpoint{4.049143in}{0.540625in}}%
\pgfpathlineto{\pgfqpoint{4.055198in}{0.540625in}}%
\pgfpathlineto{\pgfqpoint{4.061253in}{0.540625in}}%
\pgfpathlineto{\pgfqpoint{4.067308in}{0.540625in}}%
\pgfpathlineto{\pgfqpoint{4.073362in}{0.540625in}}%
\pgfpathlineto{\pgfqpoint{4.079417in}{0.540625in}}%
\pgfpathlineto{\pgfqpoint{4.085472in}{0.540625in}}%
\pgfpathlineto{\pgfqpoint{4.091527in}{0.540625in}}%
\pgfpathlineto{\pgfqpoint{4.097582in}{0.540625in}}%
\pgfpathlineto{\pgfqpoint{4.103637in}{0.540625in}}%
\pgfpathlineto{\pgfqpoint{4.109691in}{0.540625in}}%
\pgfpathlineto{\pgfqpoint{4.115746in}{0.540625in}}%
\pgfpathlineto{\pgfqpoint{4.121801in}{0.540625in}}%
\pgfpathlineto{\pgfqpoint{4.127856in}{0.540625in}}%
\pgfpathlineto{\pgfqpoint{4.133911in}{0.540625in}}%
\pgfpathlineto{\pgfqpoint{4.139966in}{0.540625in}}%
\pgfpathlineto{\pgfqpoint{4.146020in}{0.540625in}}%
\pgfpathlineto{\pgfqpoint{4.152075in}{0.540625in}}%
\pgfpathlineto{\pgfqpoint{4.158130in}{0.540625in}}%
\pgfpathlineto{\pgfqpoint{4.164185in}{0.540625in}}%
\pgfpathlineto{\pgfqpoint{4.170240in}{0.540625in}}%
\pgfpathlineto{\pgfqpoint{4.176295in}{0.540625in}}%
\pgfpathlineto{\pgfqpoint{4.182349in}{0.540625in}}%
\pgfpathlineto{\pgfqpoint{4.188404in}{0.540625in}}%
\pgfpathlineto{\pgfqpoint{4.194459in}{0.540625in}}%
\pgfpathlineto{\pgfqpoint{4.200514in}{0.540625in}}%
\pgfpathlineto{\pgfqpoint{4.206569in}{0.540625in}}%
\pgfpathlineto{\pgfqpoint{4.212624in}{0.540625in}}%
\pgfpathlineto{\pgfqpoint{4.218678in}{0.540625in}}%
\pgfpathlineto{\pgfqpoint{4.224733in}{0.540625in}}%
\pgfpathlineto{\pgfqpoint{4.230788in}{0.540625in}}%
\pgfpathlineto{\pgfqpoint{4.236843in}{0.540625in}}%
\pgfpathlineto{\pgfqpoint{4.242898in}{0.540625in}}%
\pgfpathlineto{\pgfqpoint{4.248953in}{0.540625in}}%
\pgfpathlineto{\pgfqpoint{4.255007in}{0.540625in}}%
\pgfpathlineto{\pgfqpoint{4.261062in}{0.540625in}}%
\pgfpathlineto{\pgfqpoint{4.267117in}{0.540625in}}%
\pgfpathlineto{\pgfqpoint{4.273172in}{0.540625in}}%
\pgfpathlineto{\pgfqpoint{4.279227in}{0.540625in}}%
\pgfpathlineto{\pgfqpoint{4.279227in}{2.827155in}}%
\pgfpathlineto{\pgfqpoint{4.279227in}{2.827155in}}%
\pgfpathlineto{\pgfqpoint{4.273172in}{2.839398in}}%
\pgfpathlineto{\pgfqpoint{4.267117in}{2.851477in}}%
\pgfpathlineto{\pgfqpoint{4.261062in}{2.863391in}}%
\pgfpathlineto{\pgfqpoint{4.255007in}{2.875139in}}%
\pgfpathlineto{\pgfqpoint{4.248953in}{2.886722in}}%
\pgfpathlineto{\pgfqpoint{4.242898in}{2.898139in}}%
\pgfpathlineto{\pgfqpoint{4.236843in}{2.909390in}}%
\pgfpathlineto{\pgfqpoint{4.230788in}{2.920477in}}%
\pgfpathlineto{\pgfqpoint{4.224733in}{2.931398in}}%
\pgfpathlineto{\pgfqpoint{4.218678in}{2.942155in}}%
\pgfpathlineto{\pgfqpoint{4.212624in}{2.952748in}}%
\pgfpathlineto{\pgfqpoint{4.206569in}{2.963178in}}%
\pgfpathlineto{\pgfqpoint{4.200514in}{2.973446in}}%
\pgfpathlineto{\pgfqpoint{4.194459in}{2.983553in}}%
\pgfpathlineto{\pgfqpoint{4.188404in}{2.993500in}}%
\pgfpathlineto{\pgfqpoint{4.182349in}{3.003289in}}%
\pgfpathlineto{\pgfqpoint{4.176295in}{3.012920in}}%
\pgfpathlineto{\pgfqpoint{4.170240in}{3.022396in}}%
\pgfpathlineto{\pgfqpoint{4.164185in}{3.031717in}}%
\pgfpathlineto{\pgfqpoint{4.158130in}{3.040886in}}%
\pgfpathlineto{\pgfqpoint{4.152075in}{3.049905in}}%
\pgfpathlineto{\pgfqpoint{4.146020in}{3.058775in}}%
\pgfpathlineto{\pgfqpoint{4.139966in}{3.067499in}}%
\pgfpathlineto{\pgfqpoint{4.133911in}{3.076078in}}%
\pgfpathlineto{\pgfqpoint{4.127856in}{3.084515in}}%
\pgfpathlineto{\pgfqpoint{4.121801in}{3.092813in}}%
\pgfpathlineto{\pgfqpoint{4.115746in}{3.100973in}}%
\pgfpathlineto{\pgfqpoint{4.109691in}{3.108999in}}%
\pgfpathlineto{\pgfqpoint{4.103637in}{3.116892in}}%
\pgfpathlineto{\pgfqpoint{4.097582in}{3.124655in}}%
\pgfpathlineto{\pgfqpoint{4.091527in}{3.132292in}}%
\pgfpathlineto{\pgfqpoint{4.085472in}{3.139805in}}%
\pgfpathlineto{\pgfqpoint{4.079417in}{3.147196in}}%
\pgfpathlineto{\pgfqpoint{4.073362in}{3.154469in}}%
\pgfpathlineto{\pgfqpoint{4.067308in}{3.161627in}}%
\pgfpathlineto{\pgfqpoint{4.061253in}{3.168672in}}%
\pgfpathlineto{\pgfqpoint{4.055198in}{3.175608in}}%
\pgfpathlineto{\pgfqpoint{4.049143in}{3.182438in}}%
\pgfpathlineto{\pgfqpoint{4.043088in}{3.189165in}}%
\pgfpathlineto{\pgfqpoint{4.037033in}{3.195791in}}%
\pgfpathlineto{\pgfqpoint{4.030979in}{3.202321in}}%
\pgfpathlineto{\pgfqpoint{4.024924in}{3.208757in}}%
\pgfpathlineto{\pgfqpoint{4.018869in}{3.215102in}}%
\pgfpathlineto{\pgfqpoint{4.012814in}{3.221360in}}%
\pgfpathlineto{\pgfqpoint{4.006759in}{3.227534in}}%
\pgfpathlineto{\pgfqpoint{4.000704in}{3.233627in}}%
\pgfpathlineto{\pgfqpoint{3.994650in}{3.239643in}}%
\pgfpathlineto{\pgfqpoint{3.988595in}{3.245584in}}%
\pgfpathlineto{\pgfqpoint{3.982540in}{3.251453in}}%
\pgfpathlineto{\pgfqpoint{3.976485in}{3.257254in}}%
\pgfpathlineto{\pgfqpoint{3.970430in}{3.262991in}}%
\pgfpathlineto{\pgfqpoint{3.964375in}{3.268665in}}%
\pgfpathlineto{\pgfqpoint{3.958321in}{3.274281in}}%
\pgfpathlineto{\pgfqpoint{3.952266in}{3.279840in}}%
\pgfpathlineto{\pgfqpoint{3.946211in}{3.285347in}}%
\pgfpathlineto{\pgfqpoint{3.940156in}{3.290805in}}%
\pgfpathlineto{\pgfqpoint{3.934101in}{3.296215in}}%
\pgfpathlineto{\pgfqpoint{3.928046in}{3.301581in}}%
\pgfpathlineto{\pgfqpoint{3.921992in}{3.306906in}}%
\pgfpathlineto{\pgfqpoint{3.915937in}{3.312193in}}%
\pgfpathlineto{\pgfqpoint{3.909882in}{3.317444in}}%
\pgfpathlineto{\pgfqpoint{3.903827in}{3.322663in}}%
\pgfpathlineto{\pgfqpoint{3.897772in}{3.327850in}}%
\pgfpathlineto{\pgfqpoint{3.891717in}{3.333010in}}%
\pgfpathlineto{\pgfqpoint{3.885662in}{3.338145in}}%
\pgfpathlineto{\pgfqpoint{3.879608in}{3.343256in}}%
\pgfpathlineto{\pgfqpoint{3.873553in}{3.348347in}}%
\pgfpathlineto{\pgfqpoint{3.867498in}{3.353419in}}%
\pgfpathlineto{\pgfqpoint{3.861443in}{3.358475in}}%
\pgfpathlineto{\pgfqpoint{3.855388in}{3.363516in}}%
\pgfpathlineto{\pgfqpoint{3.849333in}{3.368545in}}%
\pgfpathlineto{\pgfqpoint{3.843279in}{3.373564in}}%
\pgfpathlineto{\pgfqpoint{3.837224in}{3.378575in}}%
\pgfpathlineto{\pgfqpoint{3.831169in}{3.383578in}}%
\pgfpathlineto{\pgfqpoint{3.825114in}{3.388577in}}%
\pgfpathlineto{\pgfqpoint{3.819059in}{3.393572in}}%
\pgfpathlineto{\pgfqpoint{3.813004in}{3.398565in}}%
\pgfpathlineto{\pgfqpoint{3.806950in}{3.403557in}}%
\pgfpathlineto{\pgfqpoint{3.800895in}{3.408549in}}%
\pgfpathlineto{\pgfqpoint{3.794840in}{3.413543in}}%
\pgfpathlineto{\pgfqpoint{3.788785in}{3.418540in}}%
\pgfpathlineto{\pgfqpoint{3.782730in}{3.423541in}}%
\pgfpathlineto{\pgfqpoint{3.776675in}{3.428547in}}%
\pgfpathlineto{\pgfqpoint{3.770621in}{3.433558in}}%
\pgfpathlineto{\pgfqpoint{3.764566in}{3.438576in}}%
\pgfpathlineto{\pgfqpoint{3.758511in}{3.443600in}}%
\pgfpathlineto{\pgfqpoint{3.752456in}{3.448632in}}%
\pgfpathlineto{\pgfqpoint{3.746401in}{3.453671in}}%
\pgfpathlineto{\pgfqpoint{3.740346in}{3.458719in}}%
\pgfpathlineto{\pgfqpoint{3.734292in}{3.463776in}}%
\pgfpathlineto{\pgfqpoint{3.728237in}{3.468841in}}%
\pgfpathlineto{\pgfqpoint{3.722182in}{3.473915in}}%
\pgfpathlineto{\pgfqpoint{3.716127in}{3.478998in}}%
\pgfpathlineto{\pgfqpoint{3.710072in}{3.484089in}}%
\pgfpathlineto{\pgfqpoint{3.704017in}{3.489189in}}%
\pgfpathlineto{\pgfqpoint{3.697963in}{3.494297in}}%
\pgfpathlineto{\pgfqpoint{3.691908in}{3.499412in}}%
\pgfpathlineto{\pgfqpoint{3.685853in}{3.504536in}}%
\pgfpathlineto{\pgfqpoint{3.679798in}{3.509666in}}%
\pgfpathlineto{\pgfqpoint{3.673743in}{3.514802in}}%
\pgfpathlineto{\pgfqpoint{3.667688in}{3.519944in}}%
\pgfpathlineto{\pgfqpoint{3.661634in}{3.525091in}}%
\pgfpathlineto{\pgfqpoint{3.655579in}{3.530242in}}%
\pgfpathlineto{\pgfqpoint{3.649524in}{3.535397in}}%
\pgfpathlineto{\pgfqpoint{3.643469in}{3.540553in}}%
\pgfpathlineto{\pgfqpoint{3.637414in}{3.545711in}}%
\pgfpathlineto{\pgfqpoint{3.631359in}{3.550868in}}%
\pgfpathlineto{\pgfqpoint{3.625305in}{3.556025in}}%
\pgfpathlineto{\pgfqpoint{3.619250in}{3.561179in}}%
\pgfpathlineto{\pgfqpoint{3.613195in}{3.566330in}}%
\pgfpathlineto{\pgfqpoint{3.607140in}{3.571476in}}%
\pgfpathlineto{\pgfqpoint{3.601085in}{3.576616in}}%
\pgfpathlineto{\pgfqpoint{3.595030in}{3.581747in}}%
\pgfpathlineto{\pgfqpoint{3.588976in}{3.586870in}}%
\pgfpathlineto{\pgfqpoint{3.582921in}{3.591982in}}%
\pgfpathlineto{\pgfqpoint{3.576866in}{3.597081in}}%
\pgfpathlineto{\pgfqpoint{3.570811in}{3.602167in}}%
\pgfpathlineto{\pgfqpoint{3.564756in}{3.607236in}}%
\pgfpathlineto{\pgfqpoint{3.558701in}{3.612288in}}%
\pgfpathlineto{\pgfqpoint{3.552647in}{3.617321in}}%
\pgfpathlineto{\pgfqpoint{3.546592in}{3.622334in}}%
\pgfpathlineto{\pgfqpoint{3.540537in}{3.627323in}}%
\pgfpathlineto{\pgfqpoint{3.534482in}{3.632288in}}%
\pgfpathlineto{\pgfqpoint{3.528427in}{3.637226in}}%
\pgfpathlineto{\pgfqpoint{3.522372in}{3.642135in}}%
\pgfpathlineto{\pgfqpoint{3.516318in}{3.647015in}}%
\pgfpathlineto{\pgfqpoint{3.510263in}{3.651862in}}%
\pgfpathlineto{\pgfqpoint{3.504208in}{3.656675in}}%
\pgfpathlineto{\pgfqpoint{3.498153in}{3.661451in}}%
\pgfpathlineto{\pgfqpoint{3.492098in}{3.666190in}}%
\pgfpathlineto{\pgfqpoint{3.486043in}{3.670888in}}%
\pgfpathlineto{\pgfqpoint{3.479989in}{3.675544in}}%
\pgfpathlineto{\pgfqpoint{3.473934in}{3.680155in}}%
\pgfpathlineto{\pgfqpoint{3.467879in}{3.684720in}}%
\pgfpathlineto{\pgfqpoint{3.461824in}{3.689237in}}%
\pgfpathlineto{\pgfqpoint{3.455769in}{3.693704in}}%
\pgfpathlineto{\pgfqpoint{3.449714in}{3.698118in}}%
\pgfpathlineto{\pgfqpoint{3.443660in}{3.702477in}}%
\pgfpathlineto{\pgfqpoint{3.437605in}{3.706780in}}%
\pgfpathlineto{\pgfqpoint{3.431550in}{3.711024in}}%
\pgfpathlineto{\pgfqpoint{3.425495in}{3.715208in}}%
\pgfpathlineto{\pgfqpoint{3.419440in}{3.719329in}}%
\pgfpathlineto{\pgfqpoint{3.413385in}{3.723386in}}%
\pgfpathlineto{\pgfqpoint{3.407331in}{3.727376in}}%
\pgfpathlineto{\pgfqpoint{3.401276in}{3.731297in}}%
\pgfpathlineto{\pgfqpoint{3.395221in}{3.735148in}}%
\pgfpathlineto{\pgfqpoint{3.389166in}{3.738926in}}%
\pgfpathlineto{\pgfqpoint{3.383111in}{3.742630in}}%
\pgfpathlineto{\pgfqpoint{3.377056in}{3.746257in}}%
\pgfpathlineto{\pgfqpoint{3.371001in}{3.749806in}}%
\pgfpathlineto{\pgfqpoint{3.364947in}{3.753275in}}%
\pgfpathlineto{\pgfqpoint{3.358892in}{3.756662in}}%
\pgfpathlineto{\pgfqpoint{3.352837in}{3.759965in}}%
\pgfpathlineto{\pgfqpoint{3.346782in}{3.763183in}}%
\pgfpathlineto{\pgfqpoint{3.340727in}{3.766313in}}%
\pgfpathlineto{\pgfqpoint{3.334672in}{3.769354in}}%
\pgfpathlineto{\pgfqpoint{3.328618in}{3.772304in}}%
\pgfpathlineto{\pgfqpoint{3.322563in}{3.775161in}}%
\pgfpathlineto{\pgfqpoint{3.316508in}{3.777924in}}%
\pgfpathlineto{\pgfqpoint{3.310453in}{3.780592in}}%
\pgfpathlineto{\pgfqpoint{3.304398in}{3.783161in}}%
\pgfpathlineto{\pgfqpoint{3.298343in}{3.785632in}}%
\pgfpathlineto{\pgfqpoint{3.292289in}{3.788002in}}%
\pgfpathlineto{\pgfqpoint{3.286234in}{3.790270in}}%
\pgfpathlineto{\pgfqpoint{3.280179in}{3.792434in}}%
\pgfpathlineto{\pgfqpoint{3.274124in}{3.794493in}}%
\pgfpathlineto{\pgfqpoint{3.268069in}{3.796446in}}%
\pgfpathlineto{\pgfqpoint{3.262014in}{3.798291in}}%
\pgfpathlineto{\pgfqpoint{3.255960in}{3.800026in}}%
\pgfpathlineto{\pgfqpoint{3.249905in}{3.801652in}}%
\pgfpathlineto{\pgfqpoint{3.243850in}{3.803165in}}%
\pgfpathlineto{\pgfqpoint{3.237795in}{3.804565in}}%
\pgfpathlineto{\pgfqpoint{3.231740in}{3.805851in}}%
\pgfpathlineto{\pgfqpoint{3.225685in}{3.807022in}}%
\pgfpathlineto{\pgfqpoint{3.219631in}{3.808076in}}%
\pgfpathlineto{\pgfqpoint{3.213576in}{3.809013in}}%
\pgfpathlineto{\pgfqpoint{3.207521in}{3.809831in}}%
\pgfpathlineto{\pgfqpoint{3.201466in}{3.810529in}}%
\pgfpathlineto{\pgfqpoint{3.195411in}{3.811107in}}%
\pgfpathlineto{\pgfqpoint{3.189356in}{3.811563in}}%
\pgfpathlineto{\pgfqpoint{3.183302in}{3.811897in}}%
\pgfpathlineto{\pgfqpoint{3.177247in}{3.812108in}}%
\pgfpathlineto{\pgfqpoint{3.171192in}{3.812194in}}%
\pgfpathlineto{\pgfqpoint{3.165137in}{3.812155in}}%
\pgfpathlineto{\pgfqpoint{3.159082in}{3.811991in}}%
\pgfpathlineto{\pgfqpoint{3.153027in}{3.811700in}}%
\pgfpathlineto{\pgfqpoint{3.146973in}{3.811282in}}%
\pgfpathlineto{\pgfqpoint{3.140918in}{3.810736in}}%
\pgfpathlineto{\pgfqpoint{3.134863in}{3.810062in}}%
\pgfpathlineto{\pgfqpoint{3.128808in}{3.809259in}}%
\pgfpathlineto{\pgfqpoint{3.122753in}{3.808327in}}%
\pgfpathlineto{\pgfqpoint{3.116698in}{3.807265in}}%
\pgfpathlineto{\pgfqpoint{3.110644in}{3.806072in}}%
\pgfpathlineto{\pgfqpoint{3.104589in}{3.804748in}}%
\pgfpathlineto{\pgfqpoint{3.098534in}{3.803293in}}%
\pgfpathlineto{\pgfqpoint{3.092479in}{3.801707in}}%
\pgfpathlineto{\pgfqpoint{3.086424in}{3.799989in}}%
\pgfpathlineto{\pgfqpoint{3.080369in}{3.798138in}}%
\pgfpathlineto{\pgfqpoint{3.074315in}{3.796155in}}%
\pgfpathlineto{\pgfqpoint{3.068260in}{3.794039in}}%
\pgfpathlineto{\pgfqpoint{3.062205in}{3.791791in}}%
\pgfpathlineto{\pgfqpoint{3.056150in}{3.789409in}}%
\pgfpathlineto{\pgfqpoint{3.050095in}{3.786895in}}%
\pgfpathlineto{\pgfqpoint{3.044040in}{3.784247in}}%
\pgfpathlineto{\pgfqpoint{3.037986in}{3.781465in}}%
\pgfpathlineto{\pgfqpoint{3.031931in}{3.778551in}}%
\pgfpathlineto{\pgfqpoint{3.025876in}{3.775503in}}%
\pgfpathlineto{\pgfqpoint{3.019821in}{3.772322in}}%
\pgfpathlineto{\pgfqpoint{3.013766in}{3.769008in}}%
\pgfpathlineto{\pgfqpoint{3.007711in}{3.765560in}}%
\pgfpathlineto{\pgfqpoint{3.001657in}{3.761980in}}%
\pgfpathlineto{\pgfqpoint{2.995602in}{3.758267in}}%
\pgfpathlineto{\pgfqpoint{2.989547in}{3.754421in}}%
\pgfpathlineto{\pgfqpoint{2.983492in}{3.750442in}}%
\pgfpathlineto{\pgfqpoint{2.977437in}{3.746332in}}%
\pgfpathlineto{\pgfqpoint{2.971382in}{3.742090in}}%
\pgfpathlineto{\pgfqpoint{2.965328in}{3.737716in}}%
\pgfpathlineto{\pgfqpoint{2.959273in}{3.733210in}}%
\pgfpathlineto{\pgfqpoint{2.953218in}{3.728574in}}%
\pgfpathlineto{\pgfqpoint{2.947163in}{3.723807in}}%
\pgfpathlineto{\pgfqpoint{2.941108in}{3.718910in}}%
\pgfpathlineto{\pgfqpoint{2.935053in}{3.713884in}}%
\pgfpathlineto{\pgfqpoint{2.928999in}{3.708728in}}%
\pgfpathlineto{\pgfqpoint{2.922944in}{3.703443in}}%
\pgfpathlineto{\pgfqpoint{2.916889in}{3.698030in}}%
\pgfpathlineto{\pgfqpoint{2.910834in}{3.692489in}}%
\pgfpathlineto{\pgfqpoint{2.904779in}{3.686821in}}%
\pgfpathlineto{\pgfqpoint{2.898724in}{3.681026in}}%
\pgfpathlineto{\pgfqpoint{2.892669in}{3.675105in}}%
\pgfpathlineto{\pgfqpoint{2.886615in}{3.669059in}}%
\pgfpathlineto{\pgfqpoint{2.880560in}{3.662888in}}%
\pgfpathlineto{\pgfqpoint{2.874505in}{3.656593in}}%
\pgfpathlineto{\pgfqpoint{2.868450in}{3.650174in}}%
\pgfpathlineto{\pgfqpoint{2.862395in}{3.643633in}}%
\pgfpathlineto{\pgfqpoint{2.856340in}{3.636969in}}%
\pgfpathlineto{\pgfqpoint{2.850286in}{3.630185in}}%
\pgfpathlineto{\pgfqpoint{2.844231in}{3.623279in}}%
\pgfpathlineto{\pgfqpoint{2.838176in}{3.616254in}}%
\pgfpathlineto{\pgfqpoint{2.832121in}{3.609110in}}%
\pgfpathlineto{\pgfqpoint{2.826066in}{3.601848in}}%
\pgfpathlineto{\pgfqpoint{2.820011in}{3.594469in}}%
\pgfpathlineto{\pgfqpoint{2.813957in}{3.586973in}}%
\pgfpathlineto{\pgfqpoint{2.807902in}{3.579362in}}%
\pgfpathlineto{\pgfqpoint{2.801847in}{3.571636in}}%
\pgfpathlineto{\pgfqpoint{2.795792in}{3.563796in}}%
\pgfpathlineto{\pgfqpoint{2.789737in}{3.555843in}}%
\pgfpathlineto{\pgfqpoint{2.783682in}{3.547779in}}%
\pgfpathlineto{\pgfqpoint{2.777628in}{3.539603in}}%
\pgfpathlineto{\pgfqpoint{2.777628in}{3.539603in}}%
\pgfpathclose%
\pgfusepath{fill}%
\end{pgfscope}%
\begin{pgfscope}%
\pgfpathrectangle{\pgfqpoint{0.050000in}{0.050000in}}{\pgfqpoint{6.200000in}{3.925000in}}%
\pgfusepath{clip}%
\pgfsetrectcap%
\pgfsetroundjoin%
\pgfsetlinewidth{1.505625pt}%
\definecolor{currentstroke}{rgb}{0.156863,0.294118,0.458824}%
\pgfsetstrokecolor{currentstroke}%
\pgfsetdash{}{0pt}%
\pgfpathmoveto{\pgfqpoint{0.125610in}{0.553236in}}%
\pgfpathlineto{\pgfqpoint{0.240652in}{0.559716in}}%
\pgfpathlineto{\pgfqpoint{0.337529in}{0.567357in}}%
\pgfpathlineto{\pgfqpoint{0.422297in}{0.576181in}}%
\pgfpathlineto{\pgfqpoint{0.501010in}{0.586603in}}%
\pgfpathlineto{\pgfqpoint{0.573668in}{0.598527in}}%
\pgfpathlineto{\pgfqpoint{0.640271in}{0.611754in}}%
\pgfpathlineto{\pgfqpoint{0.700819in}{0.625984in}}%
\pgfpathlineto{\pgfqpoint{0.761367in}{0.642605in}}%
\pgfpathlineto{\pgfqpoint{0.815861in}{0.659859in}}%
\pgfpathlineto{\pgfqpoint{0.870355in}{0.679530in}}%
\pgfpathlineto{\pgfqpoint{0.918793in}{0.699244in}}%
\pgfpathlineto{\pgfqpoint{0.967232in}{0.721239in}}%
\pgfpathlineto{\pgfqpoint{1.015671in}{0.745699in}}%
\pgfpathlineto{\pgfqpoint{1.064109in}{0.772809in}}%
\pgfpathlineto{\pgfqpoint{1.112548in}{0.802755in}}%
\pgfpathlineto{\pgfqpoint{1.154932in}{0.831428in}}%
\pgfpathlineto{\pgfqpoint{1.197316in}{0.862535in}}%
\pgfpathlineto{\pgfqpoint{1.239699in}{0.896193in}}%
\pgfpathlineto{\pgfqpoint{1.282083in}{0.932513in}}%
\pgfpathlineto{\pgfqpoint{1.324467in}{0.971603in}}%
\pgfpathlineto{\pgfqpoint{1.366851in}{1.013558in}}%
\pgfpathlineto{\pgfqpoint{1.409235in}{1.058466in}}%
\pgfpathlineto{\pgfqpoint{1.451619in}{1.106402in}}%
\pgfpathlineto{\pgfqpoint{1.494003in}{1.157429in}}%
\pgfpathlineto{\pgfqpoint{1.536386in}{1.211591in}}%
\pgfpathlineto{\pgfqpoint{1.584825in}{1.277368in}}%
\pgfpathlineto{\pgfqpoint{1.633264in}{1.347291in}}%
\pgfpathlineto{\pgfqpoint{1.681702in}{1.421336in}}%
\pgfpathlineto{\pgfqpoint{1.730141in}{1.499441in}}%
\pgfpathlineto{\pgfqpoint{1.778580in}{1.581502in}}%
\pgfpathlineto{\pgfqpoint{1.833073in}{1.678362in}}%
\pgfpathlineto{\pgfqpoint{1.887567in}{1.779758in}}%
\pgfpathlineto{\pgfqpoint{1.948115in}{1.897301in}}%
\pgfpathlineto{\pgfqpoint{2.014718in}{2.031826in}}%
\pgfpathlineto{\pgfqpoint{2.087376in}{2.183711in}}%
\pgfpathlineto{\pgfqpoint{2.184254in}{2.391995in}}%
\pgfpathlineto{\pgfqpoint{2.408283in}{2.876201in}}%
\pgfpathlineto{\pgfqpoint{2.474886in}{3.013978in}}%
\pgfpathlineto{\pgfqpoint{2.535434in}{3.134120in}}%
\pgfpathlineto{\pgfqpoint{2.589928in}{3.236950in}}%
\pgfpathlineto{\pgfqpoint{2.638366in}{3.323344in}}%
\pgfpathlineto{\pgfqpoint{2.680750in}{3.394539in}}%
\pgfpathlineto{\pgfqpoint{2.723134in}{3.461195in}}%
\pgfpathlineto{\pgfqpoint{2.759463in}{3.514421in}}%
\pgfpathlineto{\pgfqpoint{2.795792in}{3.563796in}}%
\pgfpathlineto{\pgfqpoint{2.832121in}{3.609110in}}%
\pgfpathlineto{\pgfqpoint{2.862395in}{3.643633in}}%
\pgfpathlineto{\pgfqpoint{2.892669in}{3.675105in}}%
\pgfpathlineto{\pgfqpoint{2.922944in}{3.703443in}}%
\pgfpathlineto{\pgfqpoint{2.953218in}{3.728574in}}%
\pgfpathlineto{\pgfqpoint{2.983492in}{3.750442in}}%
\pgfpathlineto{\pgfqpoint{3.013766in}{3.769008in}}%
\pgfpathlineto{\pgfqpoint{3.037986in}{3.781465in}}%
\pgfpathlineto{\pgfqpoint{3.062205in}{3.791791in}}%
\pgfpathlineto{\pgfqpoint{3.086424in}{3.799989in}}%
\pgfpathlineto{\pgfqpoint{3.110644in}{3.806072in}}%
\pgfpathlineto{\pgfqpoint{3.134863in}{3.810062in}}%
\pgfpathlineto{\pgfqpoint{3.159082in}{3.811991in}}%
\pgfpathlineto{\pgfqpoint{3.183302in}{3.811897in}}%
\pgfpathlineto{\pgfqpoint{3.207521in}{3.809831in}}%
\pgfpathlineto{\pgfqpoint{3.231740in}{3.805851in}}%
\pgfpathlineto{\pgfqpoint{3.262014in}{3.798291in}}%
\pgfpathlineto{\pgfqpoint{3.292289in}{3.788002in}}%
\pgfpathlineto{\pgfqpoint{3.322563in}{3.775161in}}%
\pgfpathlineto{\pgfqpoint{3.352837in}{3.759965in}}%
\pgfpathlineto{\pgfqpoint{3.389166in}{3.738926in}}%
\pgfpathlineto{\pgfqpoint{3.425495in}{3.715208in}}%
\pgfpathlineto{\pgfqpoint{3.467879in}{3.684720in}}%
\pgfpathlineto{\pgfqpoint{3.516318in}{3.647015in}}%
\pgfpathlineto{\pgfqpoint{3.582921in}{3.591982in}}%
\pgfpathlineto{\pgfqpoint{3.843279in}{3.373564in}}%
\pgfpathlineto{\pgfqpoint{3.921992in}{3.306906in}}%
\pgfpathlineto{\pgfqpoint{3.970430in}{3.262991in}}%
\pgfpathlineto{\pgfqpoint{4.012814in}{3.221360in}}%
\pgfpathlineto{\pgfqpoint{4.049143in}{3.182438in}}%
\pgfpathlineto{\pgfqpoint{4.085472in}{3.139805in}}%
\pgfpathlineto{\pgfqpoint{4.115746in}{3.100973in}}%
\pgfpathlineto{\pgfqpoint{4.146020in}{3.058775in}}%
\pgfpathlineto{\pgfqpoint{4.176295in}{3.012920in}}%
\pgfpathlineto{\pgfqpoint{4.206569in}{2.963178in}}%
\pgfpathlineto{\pgfqpoint{4.236843in}{2.909390in}}%
\pgfpathlineto{\pgfqpoint{4.267117in}{2.851477in}}%
\pgfpathlineto{\pgfqpoint{4.303446in}{2.776549in}}%
\pgfpathlineto{\pgfqpoint{4.339775in}{2.695871in}}%
\pgfpathlineto{\pgfqpoint{4.376104in}{2.609803in}}%
\pgfpathlineto{\pgfqpoint{4.418488in}{2.503298in}}%
\pgfpathlineto{\pgfqpoint{4.466927in}{2.374925in}}%
\pgfpathlineto{\pgfqpoint{4.527475in}{2.207477in}}%
\pgfpathlineto{\pgfqpoint{4.690956in}{1.750621in}}%
\pgfpathlineto{\pgfqpoint{4.739394in}{1.623130in}}%
\pgfpathlineto{\pgfqpoint{4.781778in}{1.517201in}}%
\pgfpathlineto{\pgfqpoint{4.824162in}{1.417439in}}%
\pgfpathlineto{\pgfqpoint{4.860491in}{1.337317in}}%
\pgfpathlineto{\pgfqpoint{4.896820in}{1.262439in}}%
\pgfpathlineto{\pgfqpoint{4.933149in}{1.192932in}}%
\pgfpathlineto{\pgfqpoint{4.969478in}{1.128813in}}%
\pgfpathlineto{\pgfqpoint{4.999752in}{1.079439in}}%
\pgfpathlineto{\pgfqpoint{5.030026in}{1.033660in}}%
\pgfpathlineto{\pgfqpoint{5.060301in}{0.991349in}}%
\pgfpathlineto{\pgfqpoint{5.090575in}{0.952358in}}%
\pgfpathlineto{\pgfqpoint{5.120849in}{0.916517in}}%
\pgfpathlineto{\pgfqpoint{5.151123in}{0.883641in}}%
\pgfpathlineto{\pgfqpoint{5.181397in}{0.853540in}}%
\pgfpathlineto{\pgfqpoint{5.217726in}{0.820807in}}%
\pgfpathlineto{\pgfqpoint{5.254055in}{0.791456in}}%
\pgfpathlineto{\pgfqpoint{5.290384in}{0.765165in}}%
\pgfpathlineto{\pgfqpoint{5.326713in}{0.741627in}}%
\pgfpathlineto{\pgfqpoint{5.363042in}{0.720555in}}%
\pgfpathlineto{\pgfqpoint{5.405426in}{0.698740in}}%
\pgfpathlineto{\pgfqpoint{5.447810in}{0.679550in}}%
\pgfpathlineto{\pgfqpoint{5.490194in}{0.662655in}}%
\pgfpathlineto{\pgfqpoint{5.538633in}{0.645787in}}%
\pgfpathlineto{\pgfqpoint{5.587071in}{0.631170in}}%
\pgfpathlineto{\pgfqpoint{5.641565in}{0.617031in}}%
\pgfpathlineto{\pgfqpoint{5.702113in}{0.603754in}}%
\pgfpathlineto{\pgfqpoint{5.768716in}{0.591630in}}%
\pgfpathlineto{\pgfqpoint{5.841374in}{0.580859in}}%
\pgfpathlineto{\pgfqpoint{5.920087in}{0.571556in}}%
\pgfpathlineto{\pgfqpoint{6.004855in}{0.563753in}}%
\pgfpathlineto{\pgfqpoint{6.101732in}{0.557050in}}%
\pgfpathlineto{\pgfqpoint{6.174390in}{0.553239in}}%
\pgfpathlineto{\pgfqpoint{6.174390in}{0.553239in}}%
\pgfusepath{stroke}%
\end{pgfscope}%
\begin{pgfscope}%
\pgfpathrectangle{\pgfqpoint{0.050000in}{0.050000in}}{\pgfqpoint{6.200000in}{3.925000in}}%
\pgfusepath{clip}%
\pgfsetbuttcap%
\pgfsetroundjoin%
\pgfsetlinewidth{0.803000pt}%
\definecolor{currentstroke}{rgb}{0.000000,0.000000,0.000000}%
\pgfsetstrokecolor{currentstroke}%
\pgfsetdash{}{0pt}%
\pgfpathmoveto{\pgfqpoint{2.771951in}{0.540625in}}%
\pgfpathlineto{\pgfqpoint{2.771951in}{3.531839in}}%
\pgfusepath{stroke}%
\end{pgfscope}%
\begin{pgfscope}%
\pgfpathrectangle{\pgfqpoint{0.050000in}{0.050000in}}{\pgfqpoint{6.200000in}{3.925000in}}%
\pgfusepath{clip}%
\pgfsetbuttcap%
\pgfsetroundjoin%
\pgfsetlinewidth{0.803000pt}%
\definecolor{currentstroke}{rgb}{0.000000,0.000000,0.000000}%
\pgfsetstrokecolor{currentstroke}%
\pgfsetdash{}{0pt}%
\pgfpathmoveto{\pgfqpoint{4.284146in}{0.540625in}}%
\pgfpathlineto{\pgfqpoint{4.284146in}{2.817086in}}%
\pgfusepath{stroke}%
\end{pgfscope}%
\begin{pgfscope}%
\pgfpathrectangle{\pgfqpoint{0.050000in}{0.050000in}}{\pgfqpoint{6.200000in}{3.925000in}}%
\pgfusepath{clip}%
\pgfsetbuttcap%
\pgfsetroundjoin%
\pgfsetlinewidth{1.003750pt}%
\definecolor{currentstroke}{rgb}{0.000000,0.000000,0.000000}%
\pgfsetstrokecolor{currentstroke}%
\pgfsetdash{{1.000000pt}{1.650000pt}}{0.000000pt}%
\pgfpathmoveto{\pgfqpoint{3.906098in}{0.540625in}}%
\pgfpathlineto{\pgfqpoint{3.906098in}{3.320709in}}%
\pgfusepath{stroke}%
\end{pgfscope}%
\begin{pgfscope}%
\pgfpathrectangle{\pgfqpoint{0.050000in}{0.050000in}}{\pgfqpoint{6.200000in}{3.925000in}}%
\pgfusepath{clip}%
\pgfsetrectcap%
\pgfsetroundjoin%
\pgfsetlinewidth{0.803000pt}%
\definecolor{currentstroke}{rgb}{0.000000,0.000000,0.000000}%
\pgfsetstrokecolor{currentstroke}%
\pgfsetdash{}{0pt}%
\pgfpathmoveto{\pgfqpoint{0.050000in}{0.540625in}}%
\pgfpathlineto{\pgfqpoint{6.250000in}{0.540625in}}%
\pgfusepath{stroke}%
\end{pgfscope}%
\begin{pgfscope}%
\pgfsetbuttcap%
\pgfsetmiterjoin%
\definecolor{currentfill}{rgb}{0.000000,0.000000,0.000000}%
\pgfsetfillcolor{currentfill}%
\pgfsetlinewidth{1.003750pt}%
\definecolor{currentstroke}{rgb}{0.000000,0.000000,0.000000}%
\pgfsetstrokecolor{currentstroke}%
\pgfsetdash{}{0pt}%
\pgfsys@defobject{currentmarker}{\pgfqpoint{-0.027778in}{-0.027778in}}{\pgfqpoint{0.027778in}{0.027778in}}{%
\pgfpathmoveto{\pgfqpoint{0.027778in}{-0.000000in}}%
\pgfpathlineto{\pgfqpoint{-0.027778in}{0.027778in}}%
\pgfpathlineto{\pgfqpoint{-0.027778in}{-0.027778in}}%
\pgfpathlineto{\pgfqpoint{0.027778in}{-0.000000in}}%
\pgfpathclose%
\pgfusepath{stroke,fill}%
}%
\begin{pgfscope}%
\pgfsys@transformshift{6.174390in}{0.540625in}%
\pgfsys@useobject{currentmarker}{}%
\end{pgfscope}%
\end{pgfscope}%
\begin{pgfscope}%
\pgfsetbuttcap%
\pgfsetmiterjoin%
\definecolor{currentfill}{rgb}{0.000000,0.000000,0.000000}%
\pgfsetfillcolor{currentfill}%
\pgfsetlinewidth{1.003750pt}%
\definecolor{currentstroke}{rgb}{0.000000,0.000000,0.000000}%
\pgfsetstrokecolor{currentstroke}%
\pgfsetdash{}{0pt}%
\pgfsys@defobject{currentmarker}{\pgfqpoint{-0.027778in}{-0.027778in}}{\pgfqpoint{0.027778in}{0.027778in}}{%
\pgfpathmoveto{\pgfqpoint{-0.027778in}{0.000000in}}%
\pgfpathlineto{\pgfqpoint{0.027778in}{-0.027778in}}%
\pgfpathlineto{\pgfqpoint{0.027778in}{0.027778in}}%
\pgfpathlineto{\pgfqpoint{-0.027778in}{0.000000in}}%
\pgfpathclose%
\pgfusepath{stroke,fill}%
}%
\begin{pgfscope}%
\pgfsys@transformshift{0.125610in}{0.540625in}%
\pgfsys@useobject{currentmarker}{}%
\end{pgfscope}%
\end{pgfscope}%
\begin{pgfscope}%
\definecolor{textcolor}{rgb}{0.156863,0.294118,0.458824}%
\pgfsetstrokecolor{textcolor}%
\pgfsetfillcolor{textcolor}%
\pgftext[x=3.906098in,y=3.516959in,,base]{\color{textcolor}{\rmfamily\fontsize{10.000000}{12.000000}\selectfont\catcode`\^=\active\def^{\ifmmode\sp\else\^{}\fi}\catcode`\%=\active\def%{\%}$f(x)$}}%
\end{pgfscope}%
\begin{pgfscope}%
\definecolor{textcolor}{rgb}{0.000000,0.000000,0.000000}%
\pgfsetstrokecolor{textcolor}%
\pgfsetfillcolor{textcolor}%
\pgftext[x=3.906098in,y=0.148125in,,base]{\color{textcolor}{\rmfamily\fontsize{10.000000}{12.000000}\selectfont\catcode`\^=\active\def^{\ifmmode\sp\else\^{}\fi}\catcode`\%=\active\def%{\%}$x$}}%
\end{pgfscope}%
\begin{pgfscope}%
\definecolor{textcolor}{rgb}{0.000000,0.000000,0.000000}%
\pgfsetstrokecolor{textcolor}%
\pgfsetfillcolor{textcolor}%
\pgftext[x=3.528049in,y=0.589688in,,bottom]{\color{textcolor}{\rmfamily\fontsize{14.000000}{16.800000}\selectfont\catcode`\^=\active\def^{\ifmmode\sp\else\^{}\fi}\catcode`\%=\active\def%{\%}$\overbrace{\hspace{3.8410cm}}^{B}$}}%
\end{pgfscope}%
\end{pgfpicture}%
\makeatother%
\endgroup%

	\caption{A probability density function of a numerical random variable}\label{fig:ProbabilityDensityFunction}
\end{figure}

Describing an experiment via a random variable and its probability distribution or density is often more convenient than specifying the probability space. In fact, the probability space usually stays in the background. The question remains, however, whether there \textit{exists} a probability space \(\Omega, \cH, \Pm\) and a numerical random variable \(X\) for a given probability distribution. Fortunately, existence can be established for all practical probability models.

The concept of independence is of great importance in the study of random experiments and the construction of probability distributions. Loosely speaking, independence is about the lack of shared information between random objects.

\begin{defn}{Independent Random Variables}{}
	Let \(X\) and \(Y\) be random variables taking values in \((E, \cE)\) and \((F, \cF)\), respectively. They are said to be \textit{independent} if for any \(A \in \cE\) and \(B \in \cF\) it holds that
	\begin{equation}
		\Pm(X \in A, Y \in B) = \Pm(X \in A) \Pm(Y \in B).
		\label{eq:IndependentRandomVariablesDef}
	\end{equation}
\end{defn}

Intuitively, this means that information regarding \(X\) does not affect our knowledge of \(Y\) and vice versa. Mathematically, \Cref{eq:IndependentRandomVariablesDef} simply states that the distribution, \(\mu\) say, of \((X,Y)\) is given by the \textit{product measure} of the distributions \(\mu_X\) of \(X\) and \(\mu_Y\) of \(Y\). We can extend the independence concept to a finite or infinite collection of random variables as follows:

\begin{defn}{Independence}{}
	We say that a collection of random variables \(\{X_t, t \in \mathbb{T}\}\), with each \(X_t\) taking values in \((E_t, \cE_t)\), is independent if for any finite choice of indexes \(t_1, \dots, t_n \in \mathbb{T}\) and sets \(A_{t_1} \in \cE_{t_1}, \dots, A_{t_n} \in \cE_{t_n}\) it holds that
	\begin{equation*}
		\Pm(X_{t_1} \in A_{t_1}, \dots, X_{t_n} \in A_{t_n}) = \Pm(X_{t_1} \in A_{t_1}) \cdots \Pm(X_{t_n} \in A_{t_n}).
	\end{equation*}
\end{defn}

\begin{xattention}{Independence as a Modelling Assumption}{IndependenceAsAModellingAssumption}
	More often than not, the independence of random variables is a model \textit{assumption}, rather than a consequence. Instead of describing a random experiment via an explicit description of \(\Omega\) and \(\Pm\), we will usually model the experiment through one or more (independent) random variables
\end{xattention}

\subsection{Expectation}
The expectation of a (numerical) random variable is a weighted average of all the values that the random variable can take. The expectation of a \textit{discrete} random variable \(X\) is defined as:
\begin{equation*}
	\E X \coloneqq \sum_x x \Pm(X = x),
\end{equation*}

and for a \textit{continuous} random variable \(X\) with pdf \(f\) it is defined as:
\begin{equation*}
	\E X \coloneqq \int_{-\infty}^{\infty} x f(x) \di x.
\end{equation*}

Having to define the expectation in two ways is unsatisfactory. Also, there are cases wherein random variables are neither discrete nor continuous. How should we define the expectation for these random variables? Here is the formal definition:


\begin{defn}{Expectation}{Expectation}
	Let \((\Omega, \cH, \Pm)\) be a probability space and \(X\) a numerical random variable. The \textit{expectation} of \(X\) is defined as the integral
	\begin{equation*}
		\E X \coloneqq \int X \di \Pm.
	\end{equation*}
\end{defn}

\begin{thm}{Properties of an Expectation}{PropertiesOfAnExpectation}
	Let \(a, b \in \R\) and let \(X, Y\), and \((X_n)\) be numerical random variables for which the expectations below are well-defined. Then, the following hold:
	\begin{enumerate}
		\item (\textit{Monotonicity}): If \(X \le Y\), then \(\E X \le \E Y\).
		\item (\textit{Linearity}): \(\E(aX + bY) = a \E X + b \E Y\).
		\item (\textit{Monotone convergence}): If \(X_n \uparrow X\), then \(\E X_n \uparrow \E X\).
	\end{enumerate}
\end{thm}

\begin{lem}{Fatou}{}
	For any sequence \((X_n)\) of positive numerical random variables,
	\begin{equation*}
		\E \liminf X_n \le \liminf \E X_n.
	\end{equation*}
\end{lem}

\begin{thm}{Dominated and Bounded Convergence}{DominatedAndBoundedConvergence}
	Let \((X_n)\) be a sequence of numerical random variables for which \(\lim X_n = X\) exists. Then, the following hold:
	\begin{enumerate}
		\item (\textit{Dominated convergence}): If, for each \(n\), \(|X_n| \le Y\) for some random variable \(Y\) with \(\E Y < \infty\), then \(\E X_n \to \E X < \infty\).
		\item (\textit{Bounded convergence}): If \(|X_n| \le c\) for all \(n\) for some \(c \in \R\), then \(\E X_n \to \E X < \infty\).
	\end{enumerate}
\end{thm}

\begin{thm}{Expectation and Image Measure}{}
	Let \(X\) be a random variable taking values in \((E, \cE)\) with distribution \(\mu\) and let \(h\) be an \(\cE\)-measurable function. Then, provided the integral exists:
	\begin{equation*}
		\E h(X) = \int h(X) \di \Pm = \int_{\R} \mu(\di x) h(x) = \mu h.
	\end{equation*}
\end{thm}

\begin{thm}{Transformation Rule}{}
	Let \(X\) be a random vector with density \(f_X\) with respect to the Lebesgue measure on \((\R^n, \cB^n)\). Let \(Z \coloneqq g(\vecb{X})\), where \(g: \R^n \to \R^n\) is invertible with inverse \(g^{-1}\) and Jacobian \(|\frac{\partial \vecb{z}}{\partial \vecb{x}}|\). Then, at \(\vecb{z} = g(\vecb{x})\) the random vector \(Z\) has the following density with respect to the Lebesgue measure on \((\R^n, \cB^n)\):
	\begin{equation}
		f_Z(\vecb{z}) \coloneqq \frac{f_X(\vecb{x})}{|\frac{\partial \vecb{z}}{\partial \vecb{x}}|} = f_X(g^{-1}(\vecb{z})) \left|\frac{\partial \vecb{x}}{\partial \vecb{z}}\right|, \quad \vecb{z} \in \R^n.
		\label{eq:TransformationRuleDensity}
	\end{equation}
\end{thm}

\begin{prop}{Expectation of a Quadratic Form}{}
	Let \(\vecb{A}\) be an \(n \times n\) matrix and \(\vecb{X}\) an \(n\)-dimensional random vector with expectation vector \(\vecb{\mu}\) and covariance matrix \(\vecb{\Sigma}\). The random variable \(Y \coloneqq \vecb{X}^{\T} \vecb{A} \vecb{X}\) has expectation \(\tr(\vecb{\Sigma} \vecb{A}) + \vecb{\mu}^{\T} \vecb{A} \vecb{\mu}\).
\end{prop}

\subsection{\(L^p\) Spaces}
Let \(X\) be a numerical random variable on

\begin{defn}{\(L^p\) Space}{}
	For each \(p \in [1, \infty]\) the space \(L^p\) is comprised of all numerical random variables \(X\) for which \(\|X\|_p < \infty\).
\end{defn}

\begin{lem}{Jensen's Inequality}{JensensInequality}
	Let \(h : \mathcal{X} \to \R\) be a convex function and let \(X\) be a random variable taking values in \(\mathcal{X}\), with expectation vector \(\E X\). Then,
	\begin{equation*}
		\E h(X) \ge h(\E X).
	\end{equation*}
\end{lem}

\begin{thm}{Properties of the \(L^p\) Norm}{}
	Let \(X\) and \(Y\) be numerical random variables on \((\Omega, \cH, \Pm)\). Then, the following hold:
	\begin{enumerate}
		\item (\textit{Positivity}): \(\|X\|_p \ge 0\), and \(\|X\|_p = 0 \iff X = 0\) almost surely.
		\item (\textit{Multiplication by a constant}): \(\|c X\|_p = |c| \|X\|_p\).
		\item (\textit{Minkowski's (triangle) inequality}):
		      \begin{equation*}
			      |\|X\|_p - \|Y\|_p| \le \|X + Y\|_p \le \|X\|_p + \|Y\|_p.
		      \end{equation*}
		\item (\textit{Hölder's inequality}): For \(p, q, r \in [1, \infty]\) with \(\frac{1}{p} + \frac{1}{q} = \frac{1}{r}\),
		      \begin{equation}
			      \|XY\|_r \le \|X\|_p \|Y\|_q.
			      \label{eq:HoldersInequality}
		      \end{equation}
		\item (\textit{Monotonicity}): If \(1 \le p < q \le \infty\), then \(\|X\|_p \le \|X\|_q\).
	\end{enumerate}
\end{thm}

\subsection{Integral Transforms}
Expectations are used in many probabilistic analyses. In this section, we highlight their use in integral transforms. Many calculations and manipulations involving the use of such transforms facilitate the analysis of probability distributions. We describe two major types of transforms.

\begin{defn}{Moment Generating Function}{MomentGeneratingFunction}
	The \textit{moment generating function (MGF)} of a numerical random variable \(X\) with distribution \(\mu\) is the function \(M : \R \to [0, \infty]\), given by
	\begin{equation*}
		M(s) \coloneqq \E \e^{sX} = \int_{-\infty}^{\infty} \mu(\di x) \e^{sx}, \quad s \in \R.
	\end{equation*}
\end{defn}

\begin{thm}{Uniqueness of the MGF}{UniquenessOfTheMGF}
	A measure \(\mu\) is determined by its MGF \(M\) if the latter is finite in a neighborhood of \(0\).
\end{thm}

\begin{thm}{Taylor's Theorem for MGFs}{TaylorsTheoremForMGFs}
	Suppose \(M(s) \coloneqq \E \e^{sX}\) is finite in the neighborhood \((-s_0, s_0)\) of \(0\). Then, \(M\) has the Taylor expansion
	\begin{equation*}
		M(s) = \sum_{k=0}^{\infty} \frac{s^k}{k!} \E X^k, \quad |s| < s_0.
	\end{equation*}
	In particular, all moments \(\E X^k, k \in \N\) are finite.
\end{thm}

\begin{thm}{Product of MGFs}{}
	If \(X\) and \(Y\) are independent, then \(M_{X+Y}(s) = M_X(s) M_Y(s)\).
\end{thm}

\begin{defn}{Characteristic Function}{}
	The \textit{characteristic function} of a numerical random variable \(X\) with distribution \(\mu\) is the function \(\psi : \R \to \C\), defined by
	\begin{equation*}
		\psi(r) \coloneqq \E \e^{\mathrm{i} r X} = \E \cos(r X) + \mathrm{i} \E \sin(r X) = \int_{-\infty}^{\infty} \mu(\di x) \e^{\mathrm{i} r x}, \quad r \in \R.
	\end{equation*}
\end{defn}

\begin{thm}{Inversion and Uniqueness of the Characteristic Function}{}
	If the probability measure \(\mu\) has characteristic function \(\psi\), and if \(\mu\{a\} = \mu\{b\} = 0\), then
	\begin{equation}
		\mu((a, b]) = \lim_{t \to \infty} \frac{1}{2\pi} \int_{\R} \di r \, \I_{[-t, t]}(r) \frac{\e^{-\mathrm{i} r a} - \e^{-\mathrm{i} r b}}{\mathrm{i} r} \psi(r) \coloneqq \lim_{t \to \infty} I_t.
		\label{eq:CharacteristicFunctionInversion}
	\end{equation}
	In particular, the characteristic function determines \(\mu\) uniquely. If, in addition, \(\int_{\R} \di r \, |\psi(r)| < \infty\), then \(\mu\) has a density \(f\) with respect to the Lebesgue measure, given by
	\begin{equation}
		f(x) \coloneqq \frac{1}{2\pi} \int_{-\infty}^{\infty} \di r \, \e^{-\mathrm{i} r x} \psi(r).
		\label{eq:CharacteristicFunctionDensity}
	\end{equation}
\end{thm}

\begin{prop}{Properties of the Characteristic Function}{}
	\begin{enumerate}
		\item (\textit{Symmetry}): The random variables \(X\) and \(-X\) are identically distributed if and only if the characteristic function of \(X\) is real-valued.
		\item (\textit{Convolution}): If \(X\) and \(Y\) are independent, then \(\psi_{X+Y} = \psi_X \psi_Y\).
		\item (\textit{Affine transformation}): If \(a\) and \(b\) are real numbers, then
		      \begin{equation*}
			      \psi_{aX+b}(r) = \e^{\mathrm{i} r b} \psi_X(a r), \quad r \in \R.
		      \end{equation*}
		\item (\textit{Taylor's theorem}): If \(\E |X^n| < \infty\), then
		      \begin{equation*}
			      \psi(r) = \sum_{k=0}^n \frac{\E X^k}{k!} (\mathrm{i} r)^k + \mathcal{O}\!\left(r^n\right), \quad r \in \R.
		      \end{equation*}
	\end{enumerate}
\end{prop}

\subsection{Information and Independence}
In this section, we have another look at the role of \(\sigma\)-algebras and independence in random experiments.

\begin{defn}{\(\sigma\)-Algebra Generated by a Random Variable}{}
	Let \(X\) be a random variable taking values in \((E, \cE)\). The sets \(\{X \in A\}, A \in \cE\) form a \(\sigma\)-algebra, called the \textit{\(\sigma\)-algebra generated by \(X\)}, and we write \(\sigma X\).
\end{defn}

\begin{thm}{\(\sigma\)-Algebra Generated by a Stochastic Process}{}
	For the stochastic process \(X \coloneqq \{X_t, t \in \mathbb{T}\}\), the \(\sigma\)-algebra generated by \(X\) is the smallest \(\sigma\)-algebra on \(\Omega\) that is generated by the union of the \(\sigma\)-algebras \(\sigma X_t, t \in \mathbb{T}\). We write \(\bigvee_{t \in \mathbb{T}} \sigma X_t\) or \(\sigma \{X_t, t \in \mathbb{T}\}\).
\end{thm}

\begin{thm}{\(\sigma X\)-Numerical Random Variables}{SigmaXNumericalRandomVariables}
	Let \(X\) be a random variable taking values in a measurable space \((E, \cE)\). A mapping \(V : \Omega \to \overline{\R}\) belongs to \(\sigma X\) if and only if \(V = h \circ X\) for some \(h \in \cE\).
\end{thm}

\begin{defn}{Filtration}{}
	Let \(\mathbb{T}\) be a subset of \(\R\). A \textit{filtration} is an increasing collection \(\{\cF_t : t \in \mathbb{T}\}\) of sub-\(\sigma\)-algebras of \(\cH\); that is, \(\cF_s \subseteq \cF_t\) whenever \(s < t\).
\end{defn}

\begin{thm}{Stochastic Process with Arbitrary Index Set}{}
	For each \(t\) in an arbitrary set \(\mathbb{T}\), let \(X_t\) be a random variable taking values in a measurable set \((E_t, \cE_t)\). Then, the mapping \(V : \Omega \to \overline{\R}\) belongs to \(\sigma\{X_t, t \in \mathbb{T}\}\) if and only if there exists a sequence \((t_n)\) in \(\mathbb{T}\) and a function \(h\) in \(\bigotimes_{n} \cE_{t_n}\) such that
	\begin{equation*}
		V = h(X_{t_1}, X_{t_2}, \dots).
	\end{equation*}
\end{thm}

\begin{defn}{Independence for \(\sigma\)-Algebras}{}
	Let \(\{\cF_i, i = 1, \dots, n\}\) be a collection of sub-\(\sigma\)-algebras of \(\cH\). The \(\{\cF_i\}\) are said to be (\textit{mutually}) \textit{independent} or form an \textit{independency} if
	\begin{equation}
		\E V_1 \cdots V_n = \E V_1 \cdots \E V_n
		\label{eq:SigmaAlgebraIndependenceDef}
	\end{equation}
	for all positive random variables \(V_i \in \cF_i, i = 1, \dots, n\).
\end{defn}

\begin{thm}{Independence of \(\sigma\)-Algebras}{}
	Sub-\(\sigma\)-algebras \(\cF_1, \dots, \cF_n\) of \(\cH\) generated by p-systems \(\mathcal{C}_1, \dots, \mathcal{C}_n\) are independent if and only if for all \(H_i \in \mathcal{C}_i \cup \{\Omega\}, i = 1, \dots, n\), it holds that:
	\begin{equation}
		\Pm(H_1 \cap \dots \cap H_n) = \Pm(H_1) \cdots \Pm(H_n).
		\label{eq:ProductMeasureOnPiSystem}
	\end{equation}
\end{thm}

\begin{thm}{Independence of Random Variables}{}
	Random variables \(X\) and \(Y\) taking values in \((E, \cE)\) and \((F, \cF)\) are independent if and only if
	\begin{equation}
		\E g(X) h(Y) = \E g(X) \E h(Y), \quad g \in \cE_+, h \in \cF_+.
		\label{eq:IndependenceViaExpectation}
	\end{equation}
\end{thm}

\subsection{Important Stochastic Processes}
We conclude by highlighting some key classes of stochastic processes, including Gaussian processes, Poisson random measures, and \Levy{} processes.

\begin{defn}{Gaussian Process}{}
	A real-valued stochastic process \(\{X_t, t \in \mathbb{T}\}\) is said to be \textit{Gaussian} if all its finite-dimensional distributions are Gaussian (normal); that is, if the vector \(\bigl[X_{t_1}, \dots, X_{t_n}\bigr]^{\T}\) is multivariate normal for any choice of \(n\) and \(t_1, \dots, t_n \in \mathbb{T}\).
\end{defn}

\begin{defn}{Poisson Random Measure}{}
	A random measure \(N\) on \((E, \cE)\) is said to be a \textit{Poisson random measure} with mean measure \(\mu\) if the following properties hold:
	\begin{enumerate}
		\item \(N(A) \sim \Poi(\mu(A))\) for any set \(A \in \cE\).
		\item For any selection of disjoint sets \(A_1, \dots, A_n \in \cE\), the random variables \(N(A_1), \dots, N(A_n)\) are independent.
	\end{enumerate}
\end{defn}

\begin{prop}{Poisson Random Measure on a Product Space}{}
	Let \(X \coloneqq \{X_i, i \in I\}\) be the atoms of a Poisson random measure \(N\) on \((E, \cE)\) with mean measure \(\mu\), and let \(\{Y_i, i \in \N\}\) be \(\iid\) random variables on \((F, \cF)\) with distribution \(\pi\), independent of \(X\). Then, \(\{(X_i, Y_i), i \in I\}\) form the atoms of a Poisson random measure \(M\) on \((E \times F, \cE \otimes \cF)\) with mean measure \(\mu \otimes \pi\).
\end{prop}

\begin{thm}{Conditioning on the Total Number of Points}{}
	Let \(N\) be a Poisson random measure on \((E, \cE)\) with mean measure \(\mu\) satisfying \(\mu(E) \coloneqq c < \infty\). Then, conditional on \(N(E) = k\), the \(k\) atoms of \(N\) are \(\iid\) with probability distribution \(\mu/c\).
\end{thm}

\begin{defn}{Compound Poisson Process}{}
	The process \((X_t, t \ge 0)\) defined by
	\begin{equation*}
		X_t \coloneqq \int_{[0, t] \times \R^d} N(\di s, \di \vecb{x}) \vecb{x}, \quad t \ge 0,
	\end{equation*}
	is the \textit{compound Poisson process} corresponding to the measure \(\nu\).
\end{defn}

\begin{defn}{\Levy{} Process}{}
	The process \(X\) is said to be a \textit{\Levy{} process} if
	\begin{enumerate}
		\item the paths of \(X\) are almost surely right-continuous and left-limited, with \(X_0 = 0\), and
		\item for every \(t_1 < t_2 < t_3 < t_4\), the increments \(X_{t_2} - X_{t_1}\) and \(X_{t_4} - X_{t_3}\) are independent, and for every \(t, u \in \R_+\), the increment \(X_{t+u} - X_t\) has the same distribution as \(X_u\).
	\end{enumerate}
\end{defn}

\begin{thm}{\Levy{}-\Ito{} Decomposition}{}
	The process \(\vecb{X} \coloneqq (\vecb{X}_t, t \ge 0)\) is a \Levy{} process if and only if there exists a \(d\)-dimensional Brownian motion process \(\vecb{B} \coloneqq (\vecb{B}_t, t \ge 0)\) starting at \(\vecb{0}\) and a Poisson random measure \(N\) on \(\R_+ \times \R^d\) that is independent of \(\vecb{B}\) and has mean measure \(\text{Leb} \otimes \nu\), with \(\nu\) satisfying
	\begin{equation}
		\int_{\R^d} \nu(\di \vecb{x}) (|\vecb{x}|^2 \wedge 1) < \infty,
		\label{eq:LevyMeasureCondition}
	\end{equation}
	such that \(\vecb{X} = \vecb{B} + \vecb{U} + \vecb{Z}\), where the summands are independent, \(\vecb{U} \coloneqq (\vecb{U}_t, t \ge 0)\) is the compound Poisson process
	\begin{equation*}
		\vecb{U}_t \coloneqq \int_0^t \int_{||\vecb{x}|| > 1} N(\di s, \di \vecb{x}) \vecb{x}, \quad t \ge 0,
	\end{equation*}
	and \(\vecb{Z} \coloneqq (\vecb{Z}_t, t \ge 0)\) is the limit of compensated compound Poisson processes:
	\begin{equation*}
		\vecb{Z} \coloneqq \lim_{\delta \downarrow 0} \vecb{Z}^{\delta} \quad \text{with} \quad \vecb{Z}^{\delta}_t \coloneqq \int_0^t \int_{\delta \le ||\vecb{x}|| \le 1} \left[\vecb{x} N(\di s, \di \vecb{x}) - \di s \, \vecb{x} \nu(\di \vecb{x})\right], \quad t \ge 0.
	\end{equation*}
\end{thm}
